\documentclass[conference]{IEEEtran}
\IEEEoverridecommandlockouts

\usepackage{cite}
\usepackage{amsmath,amssymb,amsfonts}
\usepackage{mathtools}   % \coloneqq for definitional equalities (Reviewer 1)
\usepackage{graphicx}
\usepackage{booktabs}
\usepackage{array}
\usepackage{textcomp}
\usepackage{url}
\usepackage{stfloats}   % allow full-width figure* at column bottom

\graphicspath{{figures/}{../figures/}}

\begin{document}

% NOTE: title is LOCKED to the accepted submission (paper 122) for the proceedings.
\title{Label-Free, Dependency-Aware Channel Selection\\
for Coupled Multi-Channel Sensor Data}

% Author block: AUA affiliation uses the Akian College + Engineering Research Center
% naming from the ACSE instructions email; Orbeli + GWU kept for the co-author.
\author{\IEEEauthorblockN{Narek Meloyan\textsuperscript{1,2,3}, Narine Sarvazyan\textsuperscript{1,2,3,4}}
\IEEEauthorblockA{\textsuperscript{1}Zaven P.\ and Sonia Akian College of Science and Engineering, American University of Armenia, Yerevan, Armenia\\
\textsuperscript{2}Engineering Research Center, American University of Armenia, Yerevan, Armenia\\
\textsuperscript{3}L.A.\ Orbeli Institute of Physiology NAS RA, Yerevan, Armenia\\
\textsuperscript{4}George Washington University, Washington, DC, USA\\
Email: narek\_meloyan@edu.aua.am, phynas@gwu.edu}}

\maketitle

\begin{abstract}
Dimensionality reduction is a central step in high-dimensional pipelines, from
multi-channel sensor arrays and biomedical imaging to interpretable machine
learning. Without labels this usually relies on methods such as PCA or t-SNE, which
optimize label-free proxies (variance or neighborhood structure) rather than class
discriminability, and can discard a weak discriminative signal dominated by
redundant, high-variance channels. This paper presents a label-free,
dependency-aware method that keeps a subset of the original channels rather than a
projection, so the result stays interpretable. A group-structured autoencoder
reconstructs all channels; each latent factor is perturbed, and the resulting
per-channel reconstruction change gives a label-free relevance score that a
relevance-redundancy rule turns into an ordered shortlist, changing only the encoder
between modalities. Judged by a downstream classifier on the selected channels, the
method matches a supervised mutual-information selector without labels, and beats
PCA and variance ranking where the discriminative signal is not the loudest. On a
controlled benchmark, where the informative channels are also the highest-variance
ones, PCA recovers them too and the method matches the supervised selector, while
variance and random selection lag far behind. On motor-imagery
electroencephalography (EEG) the informative rhythm is weaker than the background,
so PCA and variance fall to about random, while the method stays above all of them
and matches the supervised selector using six of 22 electrodes. On biomedical
hyperspectral imaging the same method keeps or improves accuracy while removing up
to 95\% of the bands. By accurately modeling channel dependencies without data
labels, this architecture provides a powerful, interpretable alternative to
traditional dimensionality reduction in complex sensor arrays.
\end{abstract}

\begin{IEEEkeywords}
unsupervised feature selection, autoencoders, dependency-aware selection,
interpretability, brain-computer interfaces, hyperspectral imaging.
\end{IEEEkeywords}

\section{Introduction}
Complex sensors produce many channels that overlap heavily and arrive in natural
groups. For cost, power, bandwidth, and interpretability one often wants to keep a
handful of the actual channels rather than a learned mixture of them, and the
difficulty is deciding which, when no labels are available.

Unsupervised filters such as variance ranking or principal-component analysis (PCA)
score each channel by its own marginal statistics. They are fast but do not model
how channels depend on one another, and being variance-driven they fail when the
informative signal is not the loudest channel. Supervised, dependency-aware
selectors (mRMR, JMI) model conditional relevance but need labels. Few methods are
at once label-free, dependency-aware, and discrete, returning real channels rather
than a projection. The method proposed here occupies that niche; its central value
is modeling channel dependence without labels. Where the discriminative channels are
the loudest, marginal methods such as PCA recover them; where they are not, modeling
how channels depend on one another separates the informative channels from the loud,
redundant ones.

The presented approach is a label-free, dependency-aware selector that returns real
channels, built from a group-structured autoencoder with a compressive bottleneck
and perturbation-based attribution. A study of downstream classification accuracy
across four modalities shows that its channels match a supervised selector without
labels, and beat PCA and variance where the discriminative signal is not the
loudest; a simple label-using check completes the account by marking, in advance,
the redundant datasets on which selection of any kind offers little.

\section{Related Work}
Channel and feature selection methods fall into filter, wrapper, and embedded
families. Variance ranking and the Laplacian score~\cite{ref_lapscore} are
unsupervised but marginal. Dependency-aware selectors such as mRMR and
JMI~\cite{ref_mrmr,ref_jmi} model conditional relevance but are supervised. Among
unsupervised deep methods, the concrete autoencoder~\cite{ref_concrete} learns a
differentiable discrete selection layer, and in hyperspectral imaging attention
networks such as BS-Net~\cite{ref_bsnet} rank spectral bands. In EEG, channel
selection~\cite{ref_alotaiby} is mostly supervised, and wearable human-activity
recognition is also widely studied~\cite{ref_shiri}. On redundant data, random selection is a strong baseline
that many unsupervised selectors fail to beat~\cite{ref_worsethanrandom}; the check
proposed here makes that measurable.

\section{Method}
\begin{figure*}[t]
\centering
\includegraphics[width=0.55\textwidth]{fig_method.png}
\caption{The three-stage pipeline. A group-structured autoencoder is trained by
reconstruction; latent perturbations give a per-channel relevance (the
reconstruction change $I_{g,c}$ in Eq.~\eqref{eq:relevance}); and a
relevance-redundancy rule selects an ordered subset. The engine is the same across
domains; only the encoder and decoder change to match the regular axis of the data.}
\label{fig:method}
\end{figure*}
\subsection{Problem and Group Structure}
A multi-channel observation is organized into $G$ groups (a positive integer). The
grouping is the sensor's physical acquisition metadata, that is, one group per
inertial unit, excitation wavelength, scalp region, or process subsystem. It is
therefore known in advance and uses no class labels, and the method also runs with
$G{=}1$ when no grouping is available. Group $g$ contributes $C_g$ channels over a
regular axis (time or space) of length $T$, so a windowed sample is the matrix
$X_g\in\mathbb{R}^{T\times C_g}$ and the full sample the tuple $X=(X_1,\dots,X_G)$,
of which there are $N$ in the dataset. The goal is an ordered subset
$S\subseteq\{(g,c)\}$ of size $|S|=K$ (the channel budget, a positive integer) that
is informative and non-redundant, using no labels.

\subsection{Compressive Group Autoencoder}
Each group is encoded independently, the per-group features are fused by averaging,
and the code is decoded per group:
\begin{equation}
z \coloneqq \frac{1}{G}\sum_{g=1}^{G} E_g(X_g), \qquad \hat{X}_g \coloneqq D_g(z),
\end{equation}
where $E_g,D_g$ are the per-group encoder and decoder and $z\in\mathbb{R}^{d}$ is a
latent vector of size $d$. The encoder is the only modality-specific part: 1-D
temporal convolutions for signals and 2-D spatial convolutions for images. The
encoder pools the regular axis to a per-window bottleneck, so $z$ has no time or
space extent. This keeps the code from absorbing per-channel noise, which is
necessary for the attribution below to remain a reliable importance signal on
non-image data. Training minimizes reconstruction error only,
$\mathcal{L}\coloneqq\sum_{g=1}^{G}\|X_g-D_g(z)\|_2^2$, with no labels.

\subsection{Perturbation Attribution and Selection}
The $d$ latent coordinates are ranked by their variance across the $N$ samples; the
top $m$ are kept. For a retained coordinate $j\in\{1,\dots,m\}$ and a magnitude
$\delta$ from a finite set $\mathcal{D}$, the code is shifted
$z^{(j,\delta)}\coloneqq z+\delta\,e_j$ (with $e_j$ the $j$-th canonical basis vector), and
the scalar relevance of channel $c$ in group $g$ is the accumulated reconstruction
change:
\begin{equation}
I_{g,c} \coloneqq \sum_{j=1}^{m}\sum_{\delta\in\mathcal{D}} w_j\,
\frac{1}{NT}\sum_{n=1}^{N}\sum_{t=1}^{T}
\Big|\big[D_g(z^{(j,\delta)})-D_g(z)\big]_{n,t,c}\Big|,
\label{eq:relevance}
\end{equation}
averaged over the $N$ samples and $T$ axis positions, with scalar weight $w_j$
scoring coordinate $j$ by its variance rank. A channel whose reconstruction is
sensitive to the latent factors is one the joint representation depends upon.
Relevances are normalized within group,
$\tilde{I}_{g,c}\coloneqq I_{g,c}/\max_{1\le c'\le C_g} I_{g,c'}$, and channels are
chosen greedily by a relevance-redundancy rule, i.e.\ maximal marginal relevance
(Fig.~\ref{fig:method}):
\begin{equation}
c^\star = \arg\max_{(g,c)\notin S}
\Big[\tilde{I}_{g,c} - \lambda \max_{s\in S}\, \mathrm{sim}\big((g,c),s\big)\Big],
\end{equation}
with $\mathrm{sim}$ the cosine similarity between channel profiles and scalar
$\lambda\in[0,1]$ the redundancy penalty; the output is ordered and can be truncated
to any budget $K$.

\section{Experimental Setup}
Selection is always unsupervised; labels are used only by a downstream classifier
to score the selected-$K$ channels (accuracy or macro-F1) against variance,
PCA-loading, a supervised mutual-information (MI) selector, and random selection.

\emph{Controlled benchmark.} Grouped time-series with four independent latent
factors, each in a cluster of five correlated channels of decreasing amplitude,
plus 44 pure-noise channels (64 channels, 16 classes); covering all four factors is
required to classify. A $k$-nearest-neighbor classifier scores the subset.

\emph{Motor-imagery EEG.} BCI Competition IV-2a~\cite{ref_bciiv2a}: 22 electrodes in
four scalp-region groups, four classes, nine subjects, band-passed to 8 to 30\,Hz.
Because motor-imagery patterns are subject-specific, the standard within-subject
protocol is used (train session~1, test session~2) with a
common-spatial-patterns (CSP) and LDA classifier on the selected electrodes.

\emph{Biomedical imaging.} Multi-excitation fluorescence cubes (lichen tissue, 192
bands, four classes; collagen sponges, 158 bands, three classes) with a spatial
encoder, verified by per-pixel accuracy.

\section{Results and Discussion}
\subsection{Selected-Channel Accuracy in Two Regimes}
Table~\ref{tab:perf} reports the downstream classifier's score on the channels each
method selects. The label-free method stays close to the supervised MI selector in
both regimes (within $0.05$ accuracy on synthetic and $0.01$ macro-F1 on EEG)
without labels, and above variance and random. The marginal methods differ between
the two regimes. On the
synthetic benchmark the informative channels lie in the high-variance subspace, so
PCA recovers them ($0.95$, close to supervised); variance ranking, however, covers
only half the factors and stays near random ($0.24$). On EEG the informative
mu/beta rhythm is not the loudest signal, since alpha and ocular activity are
stronger, so PCA and variance do no better than random ($0.48$ and $0.47$ against
$0.49$) and select loud but uninformative electrodes, while the proposed method
recovers the informative ones. On the synthetic benchmark, eight of 64 channels
beat the full noisy set ($0.94$ against $0.87$), so selection also denoises, though
the bottleneck occasionally under-encodes a factor (mean $0.70\pm0.29$ over draws).

\begin{table}[t]
\caption{Downstream classifier score on the selected channels; no labels are used
for selection.}
\label{tab:perf}
\centering
\footnotesize
\setlength{\tabcolsep}{6pt}
\renewcommand{\arraystretch}{1.15}
\begin{tabular}{@{}lcc@{}}
\toprule
Selector & Synthetic & EEG \\
 & acc., $K{=}8$/64 & F1, $K{=}6$/22 \\
\midrule
All channels (full set) & 0.87 & 0.58 \\
Supervised MI & 0.99 & 0.53 \\
\textbf{Proposed (label-free)} & 0.94 & 0.52 \\
PCA & 0.95 & 0.48 \\
Variance & 0.24 & 0.47 \\
Random-$K$ & 0.25 & 0.49 \\
\bottomrule
\end{tabular}
\end{table}

\subsection{Motor-Imagery EEG}
Figure~\ref{fig:eeg} shows the result across budgets. At six of 22 electrodes the
label-free method reaches $0.517$ macro-F1 (about 89\% of the full $0.578$), above
PCA, variance, and random and close to the supervised MI selector. PCA is no better
than random here: it selects the loud alpha and ocular channels and misses the
sensorimotor cortex, which the reconstruction-based relevance finds. PCA catches up
only at larger budgets, and the method is at or above random for eight of the nine
subjects.

\begin{figure}[!t]
\centering
\includegraphics[width=0.74\linewidth]{fig_eeg.png}
\caption{Motor-imagery EEG (BCI-IV-2a, within-subject, CSP and LDA). At six of 22
electrodes the label-free method (red) is above PCA (purple), variance (orange), and
random (grey), and close to the supervised selector (blue). PCA is no better than
random at small budgets and catches up only at larger ones.}
\label{fig:eeg}
\end{figure}

\subsection{Biomedical Imaging and the Boundary}
On multi-excitation fluorescence cubes the same method (with a spatial encoder)
gives large reductions that act as denoising: nine of 192 bands reach 89.4\%
accuracy, a 95\% reduction and above the 88.2\% full-band value, and 30 of 158
bands raise accuracy from 79.8\% to 85.6\%. These wins share a structure: the
informative channels are neither redundant copies nor the loudest, and they are a
small fraction of the input (four of 64 synthetic factors, nine of 192 Lichens
bands), so pruning the rest is where label-free selection gains most. Because that
advantage grows with the uninformative fraction of the input, these modest
benchmarks are a conservative test for a method aimed at much larger channel arrays.
Selection offers little only at the opposite extreme, where channels are so redundant that
every subset carries the same information, as on wearable human-activity recognition
(PAMAP2, Opportunity) and adjacent-band remote-sensing imagery (Indian Pines): there
no unsupervised selector, the proposed one or PCA, separates from a random subset,
and a supervised selector is needed~\cite{ref_worsethanrandom}. A label-using check
(best $K$ versus random $K$ under a supervised score) flags such cases before any
modeling.

\section{Conclusion}
This paper presented a label-free, dependency-aware channel selector: a compressive
group autoencoder with perturbation attribution and relevance-redundancy selection.
Measured by downstream classification accuracy, its channels match a supervised
selector without labels and beat PCA and variance where the discriminative signal is
not the loudest, and remove up to 95\% of the bands on biomedical imaging without
loss. A simple label-using check flags, in advance, the strongly redundant datasets
on which every subset performs alike. Future work will tighten the bottleneck and
study attribution less sensitive to channel variance.

\section*{Acknowledgment}
Financial support from the European Union NAS SAR-101087403 and the United States
R44 HL120511 awards is gratefully acknowledged. Narek Chilingaryan and Kristina
Ghahramanyan are thanked for providing labelled datasets that included lichen
samples and collagen sponges.

\begin{thebibliography}{00}
\bibitem{ref_concrete} M. F. Bal\i{}n, A. Abid, and J. Zou, ``Concrete
autoencoders: Differentiable feature selection and reconstruction,'' in
\emph{Proc. ICML}, 2019.
\bibitem{ref_mrmr} H. Peng, F. Long, and C. Ding, ``Feature selection based on
mutual information,'' \emph{IEEE Trans. Pattern Anal. Mach. Intell.}, vol.~27,
no.~8, pp.~1226--1238, 2005.
\bibitem{ref_jmi} G. Brown, A. Pocock, M.-J. Zhao, and M. Luj\'an, ``Conditional
likelihood maximisation: A unifying framework for information theoretic feature
selection,'' \emph{J. Mach. Learn. Res.}, vol.~13, pp.~27--66, 2012.
\bibitem{ref_lapscore} X. He, D. Cai, and P. Niyogi, ``Laplacian score for
feature selection,'' in \emph{Proc. NeurIPS}, 2005.
\bibitem{ref_bsnet} Y. Cai, X. Liu, and Z. Cai, ``BS-Nets: An end-to-end
framework for band selection of hyperspectral image,'' \emph{IEEE Trans. Geosci.
Remote Sens.}, vol.~58, no.~3, pp.~1969--1984, 2020.
\bibitem{ref_alotaiby} T. N. Alotaiby, F. E. Abd El-Samie, S. A. Alshebeili, and
I. Ahmad, ``A review of channel selection algorithms for EEG signal processing,''
\emph{EURASIP J. Adv. Signal Process.}, vol.~2015, art.~66, 2015.
\bibitem{ref_shiri} F. M. Shiri, T. Perumal, N. Mustapha, and R. Mohamed, ``Deep
learning and federated learning in human activity recognition with sensor data: A
comprehensive review,'' \emph{Comput. Model. Eng. Sci.}, vol.~145, no.~2,
pp.~1389--1478, 2025.
\bibitem{ref_worsethanrandom} M. Rajabinasab, M. E. Houle, O. Chelly, and A.
Zimek, ``Worse than random? The importance of a baseline for unsupervised feature
selection,'' \emph{arXiv:2605.22973}, 2026.
\bibitem{ref_bciiv2a} M. Tangermann \emph{et al.}, ``Review of the BCI Competition
IV,'' \emph{Front. Neurosci.}, vol.~6, art.~55, 2012.
\end{thebibliography}

\end{document}
